\documentclass{article}
\usepackage{amsmath}
\usepackage{parskip, parselines, geometry}
\begin{document}

We starten met de evenwichtsvoorwaarde van de wet van Henry voor het oplossen van $\text{CO}_2$:
\begin{equation}
    [\text{H}_2\text{CO}_3] = K_H \cdot P_{\text{CO}_2}
\end{equation}

Vervolgens bekijken we de eerste zuurconstante $K_{a1}$:
\begin{equation}
    K_{a1} = \frac{[\text{H}^+][\text{HCO}_3^-]}{[\text{H}_2\text{CO}_3]}
\end{equation}

Omdat het zwakke zuur in een $1:1$ verhouding dissocieert, geldt in evenwicht:
\begin{equation}
    [\text{H}^+] \approx [\text{HCO}_3^-]
\end{equation}

Door vergelijking (1) en (3) te substitueren in vergelijking (2), verkrijgen we:
\begin{equation}
    K_{a1} = \frac{[\text{H}^+]^2}{K_H \cdot P_{\text{CO}_2}}
\end{equation}

Herschrijf dit om de concentratie waterstofionen kwadratisch te isoleren:
\begin{equation}
    [\text{H}^+]^2 = K_{a1} \cdot K_H \cdot P_{\text{CO}_2}
\end{equation}

Neem nu de negatieve logaritme ($-\log$) van beide leden om de overstap naar de pH-schaal te maken:
\begin{align*}
    -\log\left([\text{H}^+]^2\right) &= -\log(K_{a1} \cdot K_H \cdot P_{\text{CO}_2}) \\
    -2\log([\text{H}^+]) &= -\log(K_{a1} \cdot K_H \cdot P_{\text{CO}_2}) 
\end{align*}

Aangezien de definitie van pH gelijk is aan $-\log([\text{H}^+])$, kunnen we dit direct substitueren:
\begin{align*}
    2\text{pH} &= -\log(K_{a1} \cdot K_H \cdot P_{\text{CO}_2}) \\
    \text{pH} &= -\frac{1}{2}\log(K_{a1} \cdot K_H \cdot P_{\text{CO}_2})
\end{align*}

\end{document}